\chapter{Limitations of the Michal--Bastiani calculus}
\label{chap:limitations}

The Michal--Bastiani calculus is defined on arbitrary locally convex spaces, needs no
completeness, and its smooth maps are continuous. The price is paid elsewhere: spaces of
operators and of smooth maps carry several natural topologies, none of which makes evaluation
continuous in infinite dimensions, and as a consequence the category of smooth maps is
\emph{not} cartesian closed. This chapter collects these caveats, explains why they are not a
matter of completeness, and records the convention that we follow throughout to avoid them:
operator-valued data are always required to be smooth \emph{jointly} (in uncurried form), never
as maps into a space of operators. Nothing in this chapter is formalized; it is here to explain
the design.

Throughout, $E$ and $F$ are locally convex spaces over $\mathbb{K} \in \{\R, \mathbb{C}\}$, and
$\Lin(E, F)$ is the space of continuous linear maps.

\section{Topologies on spaces of operators}

Three topologies on $\Lin(E, F)$ will appear, each one the topology of uniform convergence on
a family of subsets of $E$:
\begin{itemize}
  \item $\Lin(E, F)_s$: uniform convergence on finite sets (pointwise convergence); on normed
    spaces this is the \emph{strong operator topology};
  \item $\Lin(E, F)_c$: uniform convergence on compact sets (the compact-open topology);
  \item $\Lin(E, F)_b$: uniform convergence on bounded sets; on normed spaces this is the
    \emph{operator norm} topology, since it suffices to look at the unit ball.
\end{itemize}
Since finite sets are compact and compact sets are bounded, $s \subseteq c \subseteq b$.

A word of caution on terminology: in operator theory ``strong topology'' means $s$, while in
the theory of locally convex spaces the ``strong topology'' (on a dual space, say) is $b$.

\begin{remark}[$s$ and $c$ agree on equicontinuous sets]
  \label{rem:s-c-equicontinuous}
  On an equicontinuous subset $H \subseteq \Lin(E, F)$ (for normed $E$, $F$: a set bounded in
  operator norm) the topologies $s$ and $c$ coincide. Indeed, given a compact $K \subseteq E$,
  a continuous seminorm $q$ on $F$ and $\varepsilon > 0$, equicontinuity gives a
  $0$-neighborhood $V$ with $q(T v) < \varepsilon$ for $T \in H$, $v \in V$; covering $K$ by
  finitely many translates $x_i + V$ reduces uniform convergence on $K$ to convergence at the
  finitely many points $x_i$. In particular, for a Banach space $E$, a sequence in
  $\Lin(E, F)$ which converges in $s$ is bounded in operator norm (uniform boundedness
  principle), hence converges in $c$: on Banach spaces the two topologies can only be
  distinguished by nets.
\end{remark}

\begin{remark}[$c$ is strictly finer than $s$ in infinite dimensions]
  \label{rem:c-finer-than-s}
  If $E$ is infinite-dimensional and $F \ne 0$, then $s \ne c$ on $\Lin(E, F)$. It suffices
  to treat $F = \mathbb{K}$, i.e.\ the dual $E'$. Every $s$-neighborhood of $0$ in $E'$
  contains the annihilator of a finite set $\{x_1, \dots, x_k\}$. If the polar $K^\circ$ of a
  compact set $K$ contained such an annihilator, then every $f \in E'$ vanishing on the
  $x_j$ would vanish on $K$ (scale $f$), and by the Hahn--Banach theorem $K$ would lie in the
  finite-dimensional span of the $x_j$. A compact set not contained in any finite-dimensional
  subspace exists (for normed $E$: $\{0\} \cup \{x_n / n\}$ with $x_n$ linearly independent
  unit vectors), so $K^\circ$ is a $c$-neighborhood of $0$ which is not an $s$-neighborhood.
\end{remark}

\begin{example}
  \label{ex:s-not-c}
  \begin{enumerate}
    \item (Incomplete $E$, sequences.) Let $E = c_{00}$ be the space of finitely supported
      sequences with the supremum norm, and $\varphi_n = n\, e_n^*$. For each $x \in E$,
      $\varphi_n(x) = 0$ for large $n$, so $\varphi_n \to 0$ in $s$. On the compact set
      $K = \{0\} \cup \{e_m / m : m \ge 1\}$ we have $\varphi_n(e_n / n) = 1$, so
      $\varphi_n \not\to 0$ in $c$. (By Remark~\ref{rem:s-c-equicontinuous}, such a sequence
      cannot exist on a Banach space.)
    \item ($E = \ell^2$, nets.) Let $A = \{\sqrt{n}\, e_n^* : n \ge 1\} \subseteq (\ell^2)'$.
      Then $0$ lies in the $s$-closure of $A$: otherwise there would be
      $x^{(1)}, \dots, x^{(k)} \in \ell^2$ and $\varepsilon > 0$ with
      $\max_j \sqrt{n}\,|x^{(j)}_n| \ge \varepsilon$ for all $n$, whence
      $\sum_j |x^{(j)}_n|^2 \ge \varepsilon^2 / n$, and summing over $n$ gives a contradiction.
      But $0$ is not in the $c$-closure of $A$: on the compact set
      $K = \{0\} \cup \{e_m / \sqrt{m}\}$ each element of $A$ takes the value $1$.
  \end{enumerate}
\end{example}

For a Banach space $E$, the topology $c$ on $E'$ is the \emph{bounded weak-$*$ topology}: the
finest topology which agrees with the weak-$*$ topology on bounded sets (Banach--Dieudonné
theorem; for the statement on convex sets see \cite[Thm.~4.44]{FHHMPZ01}).

\section{The exponential law}

For locally convex subsets $U \subseteq E_1$ and $V \subseteq E_2$ with dense interior, endow
$C^\infty(V, F)$ with the compact-open $C^\infty$-topology, and write
$f^\vee(x) = f(x, \cdot)$.

\begin{theorem}[{\cite[Thm.~1.7.30, Cor.~1.7.35]{GN26}}]
  \label{thm:exponential-law-smooth}
  For all locally convex $E_1$, $E_2$, $F$ and every $f \in C^\infty(U \times V, F)$, the map
  $f^\vee \colon U \to C^\infty(V, F)$ is smooth, and
  $\Phi \colon C^\infty(U \times V, F) \to C^\infty(U, C^\infty(V, F))$, $f \mapsto f^\vee$, is
  linear and a topological embedding. If moreover $V$ is locally compact, or
  $U \times V \times E_1 \times E_2$ is a $k_{\R}$-space, then $\Phi$ is an isomorphism of
  topological vector spaces.
\end{theorem}

\begin{remark}[Completeness plays no role]
  \label{rem:exponential-law-completeness}
  No completeness is assumed in Theorem~\ref{thm:exponential-law-smooth}. The hypothesis for the
  converse direction is topological: a $k_{\R}$-space is a Hausdorff space on which a real
  function is continuous as soon as its restrictions to compact sets are
  \cite[Rem.~A.6.15]{GN26}. It holds when $E_1$ and $E_2$ are metrizable (normed spaces,
  complete or not; Fréchet spaces), when $V$ is finite-dimensional, and also for
  $k_\omega$-spaces such as countable direct limits of finite-dimensional spaces, e.g.\
  $\R^{(\N)}$ \cite[Rem.~4.1.22, Prop.~A.7.20]{GN26}. So the absence of a countable basis of
  $0$-neighborhoods only threatens the converse direction, and only for spaces which are not
  $k_{\R}$-spaces.
\end{remark}

\section{Evaluation and cartesian closedness}

The category would be cartesian closed if $C^\infty(V, F)$ were an object and the evaluation
$\varepsilon \colon C^\infty(V, F) \times V \to F$ were smooth. This holds when $V$ is
finite-dimensional: then $\varepsilon$ is $C^{\infty, \ell}$ \cite[Lemma~1.7.25]{GN26}. It
fails in every infinite-dimensional case, and already at the linear level.

\begin{proposition}[Evaluation on the compact-open dual]
  \label{prop:evaluation-discontinuous}
  If $E$ is infinite-dimensional, the evaluation $E'_c \times E \to \mathbb{K}$,
  $(f, x) \mapsto f(x)$, is not continuous. Consequently, for every open
  $U \subseteq E$, the evaluation $C^\infty(U, \mathbb{K}) \times U \to \mathbb{K}$ is not
  continuous for the compact-open $C^\infty$-topology.
\end{proposition}

\begin{proof}
  Suppose the evaluation is continuous at $(0, 0)$. Then there are a compact $K \subseteq E$,
  $\delta > 0$ and an absolutely convex $0$-neighborhood $W \subseteq E$ with $|f(x)| \le 1$
  whenever $\sup_K |f| \le \delta$ and $x \in W$. By homogeneity, $K^\circ \subseteq
  \delta^{-1} W^\circ$, and by the bipolar theorem
  $W \subseteq \overline{W} = W^{\circ\circ} \subseteq \delta^{-1} K^{\circ\circ}$, where
  $K^{\circ\circ}$ is the closed absolutely convex hull of $K$. The absolutely convex hull of a
  precompact set is precompact, and so is its closure; hence $E$ has a precompact
  $0$-neighborhood and is finite-dimensional.

  For the second statement, the compact-open $C^\infty$-topology induces on
  $E' \subseteq C^\infty(U, \mathbb{K})$ the topology of $E'_c$ (for a linear map all
  derivatives of order $\ge 1$ are $f$ itself or zero), and continuity of the evaluation at
  $(0, x_0)$ would give continuity of $(f, h) \mapsto f(x_0 + h) - f(x_0) = f(h)$ at $(0, 0)$.
\end{proof}

The proposition covers Banach and Hilbert spaces: there the evaluation
$\Lin(E, F)_b \times E \to F$ \emph{is} continuous \cite[Rem.~2.3.9]{GN26}, but $b$ is not
the topology which the compact-open $C^\infty$-topology induces on operators. For non-normable
$E$ there is no way out at all: the evaluation $\Lin(E) \times E \to E$ is discontinuous for
\emph{every} vector topology on $\Lin(E)$ (Maissen \cite[Satz 1]{Mai63}, see also \cite[p.~2]{KM97}
and the notes to Ch.~3 of \cite{GN26}).

\begin{remark}[The convenient calculus]
  \label{rem:convenient}
  The convenient calculus of Frölicher, Kriegl and Michor \cite{KM97} is cartesian closed. It
  works on Mackey complete ($c^\infty$-complete) spaces, and this is where a completeness
  hypothesis enters. The price: the exponential law holds only as an isomorphism of
  bornological spaces (already for $E_1 = E_2 = \R$, the map $\Phi$ is discontinuous for the
  topologies of the convenient setting, see the notes to Ch.~1 of \cite{GN26}), and convenient
  smooth maps need not be continuous. On Fréchet spaces both calculi have the same smooth maps.
\end{remark}

\section{The derivative as an operator-valued map}

Let $f \colon U \to F$ be $\Ck{k}$ with $k \ge 1$, $U \subseteq E$ open.

\begin{itemize}
  \item \emph{Into $\Lin(E, F)_c$, no hypotheses.} The map $x \mapsto \dd f(x) \in
    \Lin(E, F)_c$ is $\Ck{k-1}$, for arbitrary locally convex $E$, $F$: apply the first part of
    Theorem~\ref{thm:exponential-law-smooth} to $\dd f$ and corestrict to the closed subspace
    $\Lin(E, F) \subseteq C(E, F)$ (as in the proof of \cite[Prop.~2.3.17(b)]{GN26}). Dually,
    for fixed $x$ the map $C^\infty(U, F) \to \Lin(E, F)_c$, $f \mapsto \dd f(x)$, is
    continuous linear, directly from the definition of the compact-open $C^\infty$-topology.
  \item \emph{Into $\Lin(E, F)_b$, normed spaces.} If $E$, $F$ are normed, a Michal--Bastiani
    $\Ck{k+1}$ map is Fréchet $\Ck{k}$ (Theorem~\ref{thm:mb-to-frechet}): continuity of
    $x \mapsto \dd f(x)$ in operator norm costs one derivative, and nothing for smooth maps
    (Corollary~\ref{cor:smooth-iff}). On the other hand, $f \mapsto \dd f(x)$ is \emph{not}
    continuous from $C^\infty(U, F)$ to $\Lin(E, F)_b$ when $E$ is infinite-dimensional, since
    on $E'$ the compact-open $C^\infty$-topology is $c \ne b$.
  \item \emph{Non-normable spaces.} By Maissen's theorem, there is no vector topology on
    $\Lin(E)$ making evaluation continuous, so no choice of topology turns the derivative into
    a well-behaved operator-valued map while keeping the chain rule continuous.
\end{itemize}

\section{The convention adopted}

To stay clear of all of the above, operator-valued data are always required to be smooth in
uncurried form. A smooth family of projectors on an embedded submanifold $X \subseteq F$ is a
map $E \colon X \times F \to F$, linear in the second variable, such that $(p, w) \mapsto
E_p w$ is smooth in charts (\texttt{LocallyConvex.Manifolds.ProjectorField}); the Christoffel symbols of a
connection are smooth as functions $(x, a, b) \mapsto \Gamma_x(a, b)$ of three variables
(\texttt{LocallyConvex.Manifolds.Connection}). No topology on $\Lin(F)$ is ever needed. On Banach spaces
this loses nothing: by Corollary~\ref{cor:smooth-iff}, $(x, w) \mapsto E_x w$ is then Fréchet
smooth, hence so is $x \mapsto E_x$ into $\Lin(F)$ with the operator norm; the uncurried and
the operator-norm notions of smoothness agree for $C^\infty$ data.

The convention does have visible consequences on non-normable models: the dual bundle
$T^*M = \bigcup_x (T_x M)'$ carries in general no canonical vector bundle structure
\cite[notes to Ch.~3]{GN26}, so $1$-forms are treated as smooth maps on $TM$ which are linear on
the fibres; and the adjoint action of a Lie group is handled as the smooth map
$G \times \mathfrak{g} \to \mathfrak{g}$, not as a map into $\Lin(\mathfrak{g})$.
